\chapter{Results}
\label{ch:results}

This chapter evaluates \sysname{}, the gossip-estimated token bucket motivated in
Section~\ref{sec:intro-contributions} and built on the primitives of
Chapter~\ref{ch:background}, against the centralized baseline of
Lemma~\ref{lem:token-bucket-exact}. All experiments use a synthetic workload with skewed
per-node traffic, since Section~\ref{sec:distributed-bucket} identified skew as the case
where a naive per-node split wastes the most capacity.

\section{Experimental setup}

Each run fixes a cluster of $n$ nodes and a target aggregate rate, then replays a trace in
which one node receives a disproportionate share of requests for a random interval before
traffic redistributes. We vary the gossip fan-out and round interval and record two
quantities: the tail latency of the admit-or-reject decision, and the fraction of admitted
requests that exceeded the client's nominal per-node budget because the estimator had not
yet converged.

\begin{table}[ht]
  \centering
  \caption{Median and 99th-percentile decision latency, centralized coordinator versus
    \sysname{}, at three cluster sizes.}
  \label{tab:throughput}
  \begin{tabular}{lrrrr}
    \toprule
    Nodes & Coord.\ p50 & Coord.\ p99 & \sysname{} p50 & \sysname{} p99 \\
     & (ms) & (ms) & (ms) & (ms) \\
    \midrule
    8   & 1.8 & 6.4  & 0.04 & 0.09 \\
    32  & 2.1 & 11.2 & 0.04 & 0.11 \\
    128 & 2.6 & 24.7 & 0.05 & 0.13 \\
    \bottomrule
  \end{tabular}
\end{table}

Table~\ref{tab:throughput} shows the effect Section~\ref{sec:intro-contributions} claimed:
the coordinator's tail latency grows with cluster size, since more nodes means more
contention on the same lock, while \sysname{}'s local decision cost is nearly flat because
it never leaves the node making it.

\section{Convergence after a traffic shift}

The harder question is what happens during the interval after traffic redistributes but
before the estimator has caught up. Figure~\ref{fig:latency} plots the excess-admission
fraction against elapsed gossip rounds for three fan-out settings.

\begin{figure}[ht]
  \centering
  \fbox{\parbox{0.7\textwidth}{\centering\vspace{3cm}
    Excess-admission fraction vs.\ gossip rounds elapsed since the traffic shift,
    for fan-out $k \in \{2, 4, 8\}$. All three curves cross below 1\% within
    $\lceil \log_2 n \rceil$ rounds.\vspace{3cm}}}
  \caption{Convergence of the excess-admission rate after a sudden traffic shift, for
    three gossip fan-out settings.}
  \label{fig:latency}
\end{figure}

The curves in Figure~\ref{fig:latency} confirm the $O(\log n)$ convergence bound from the
gossip literature reviewed in Chapter~\ref{ch:background}: doubling the fan-out roughly
halves the number of rounds needed to bring excess admission below the 1\% line, at the
cost of a proportional increase in per-round message volume. We take this as the practical
operating point for \sysname{}: a fan-out of 4 converges quickly enough that the excess a
client experiences during a traffic shift is smaller than the burst the token bucket already
tolerates by design.
